Table~\ref{tab:main} gives seed means $\pm$ standard deviations for the sparse family, Table~\ref{tab:robust} gives medians, relative errors and failure counts, and Table~\ref{tab:perseed} lists every seed at 64 nodes, with the dense error at 8 nodes for the degree-shift comparison. Figure~\ref{fig:mae} shows all seeds at all sizes. On the registered primary metric, mean S64 MAE (Table~\ref{tab:main}), the sum MPNN trained on the final distance is lowest ($0.23 \pm 0.19$), followed by the max MPNN trained on the final distance ($0.91 \pm 1.22$), the MLP, max+steps and sum+steps. These means are dominated by a minority of seeds (sum+steps seed 4 has S64 MAE of order $10^5$, Table~\ref{tab:perseed}), so we also report medians and per-seed failure counts; that choice was made after seeing the data.

\begin{table*}[t]\centering\caption{Sparse family, MAE (mean $\pm$ std over seeds; last column = number of seeds); bold best, underline second.}\label{tab:main}
\small\begin{tabular}{lccccc}
\toprule
Method & mae\_sparse\_n8 $\downarrow$ & mae\_sparse\_n16 $\downarrow$ & mae\_sparse\_n32 $\downarrow$ & mae\_sparse\_n64 $\downarrow$ & $n$ \\
\midrule
MPNN-sum + steps & 0.04 $\pm$ 0.01 & 0.12 $\pm$ 0.04 & 0.43 $\pm$ 0.24 & 54707.72 $\pm$ 122329.15 & 5 \\
MPNN-max & \textbf{0.01 $\pm$ 0.00} & \textbf{0.03 $\pm$ 0.01} & 0.25 $\pm$ 0.46 & \underline{0.91 $\pm$ 1.22} & 5 \\
MPNN-sum & \underline{0.02 $\pm$ 0.00} & \underline{0.03 $\pm$ 0.01} & \textbf{0.08 $\pm$ 0.06} & \textbf{0.23 $\pm$ 0.19} & 5 \\
MLP (flat adjacency) & 0.24 $\pm$ 0.00 & 0.67 $\pm$ 0.01 & 1.02 $\pm$ 0.02 & 1.33 $\pm$ 0.02 & 5 \\
\textbf{MPNN-max + steps} & 0.03 $\pm$ 0.01 & 0.08 $\pm$ 0.04 & \underline{0.24 $\pm$ 0.18} & 1.96 $\pm$ 2.05 & 5 \\
\bottomrule
\end{tabular}

\end{table*}

\begin{table*}[t]\centering\caption{Medians over 5 seeds. S/D: sparse/dense family, number = test nodes; MAE except rS64/rD64 (relative error). ``fail'' = seeds with MAE $>1$ or non-finite; div.\ = non-finite.}\label{tab:robust}
\small\begin{tabular}{lcccccccccc}\toprule
 System & S8 & S16 & S32 & S64 & D8 & D64 & rS64 & rD64 & fail S64 & fail D64 \\\midrule
MLP (flat adjacency) & 0.239 & 0.667 & 1.020 & 1.328 & 0.244 & 0.562 & 0.943 & 0.926 & 5/5 & 0/5 \\
MPNN-sum & 0.018 & 0.028 & 0.063 & 0.140 & 0.017 & 0.464 & 0.097 & 0.763 & 0/5 & 1/5 \\
MPNN-sum + steps & 0.036 & 0.108 & 0.400 & 0.460 & 0.037 & div. & 0.324 & div. & 1/5 & 5/5 \\
MPNN-max & 0.006 & 0.026 & 0.053 & 0.073 & 0.006 & 0.076 & 0.051 & 0.127 & 2/5 & 2/5 \\
MPNN-max + steps & 0.021 & 0.057 & 0.143 & 1.765 & 0.021 & 3.611 & 1.223 & 6.055 & 3/5 & 3/5 \\
\bottomrule\end{tabular}

\end{table*}

\begin{table*}[t]\centering\caption{Per-seed values (seeds 0--4): sparse MAE at 64 nodes, dense MAE at 64 nodes, and dense MAE at 8 nodes (H4 compares the last two). div.\ = non-finite; values of 100 or more are given to one significant figure ($3\mathrm{e}5=3\times10^5$).}\label{tab:perseed}
\setlength{\tabcolsep}{4pt}\small\begin{tabular}{l|ccccc|ccccc|ccccc}\toprule
 & \multicolumn{5}{c|}{S64 MAE, seed} & \multicolumn{5}{c|}{D64 MAE, seed} & \multicolumn{5}{c}{D64/D8, seed} \\
 System & 0 & 1 & 2 & 3 & 4 & 0 & 1 & 2 & 3 & 4 & 0 & 1 & 2 & 3 & 4 \\\midrule
MLP (flat adjacency) & 1.32 & 1.33 & 1.32 & 1.36 & 1.33 & 0.57 & 0.55 & 0.59 & 0.54 & 0.56 & 2.30 & 2.22 & 2.45 & 2.20 & 2.32 \\
MPNN-sum & 0.42 & 0.45 & 0.063 & 0.14 & 0.054 & 0.46 & 0.19 & 0.58 & 0.22 & div. & 24.26 & 11.53 & 40.66 & 8.39 & div. \\
MPNN-sum + steps & 0.33 & 0.46 & 0.39 & 0.67 & 3e5 & div. & div. & div. & 1e17 & div. & div. & div. & div. & 3e18 & div. \\
MPNN-max & 2.63 & 1.79 & 0.011 & 0.073 & 0.042 & 3.71 & 2.45 & 0.008 & 0.076 & 0.046 & 4e2 & 3e2 & 2.02 & 12.28 & 8.02 \\
MPNN-max + steps & 0.21 & 5.15 & 2.50 & 1.76 & 0.15 & 0.16 & 35.55 & 4.71 & 3.61 & 0.13 & 5.72 & 8e2 & 2e2 & 2e2 & 6.95 \\
\bottomrule\end{tabular}

\end{table*}

\begin{figure}[t]\centering\includegraphics[width=\columnwidth]{mae_vs_n.pdf}
\caption{MAE versus test size (log axes), median line and one dot per seed; left: sparse family, right: dense family. The axis is clipped: non-finite values and finite values of $10^5$ or more (up to order $10^{17}$, Table~\ref{tab:perseed}) are all drawn at $10^5$ (dotted line).}\label{fig:mae}\end{figure}

\begin{table}[t]\centering\caption{Tests against MPNN-max + steps (5 seeds, seed-level MAE; ``Ref.'' and ``Other'' are the group means of max+steps and of the compared system). Rows marked $\ast$ were not registered and are post hoc, uncorrected.}\label{tab:hyp}
\setlength{\tabcolsep}{4pt}\small\begin{tabular}{lllcccc}\toprule
 & System & Cell & Ref. & Other & Welch $p$ & paired $p$ \\\midrule
H1 & sum+steps & S64 & 1.955 & 5.5e+04 & 0.374 & 0.374 \\
H1 & sum+steps & S16$^\ast$ & 0.078 & 0.119 & 0.160 & 0.231 \\
H2 & max & S64 & 1.955 & 0.909 & 0.363 & 0.363 \\
H2 & max & S16$^\ast$ & 0.078 & 0.027 & 0.044 & 0.041 \\
H3 & MLP & S16 & 0.078 & 0.671 & 7.5e-07 & 4.2e-06 \\
H3 & MLP & S64 & 1.955 & 1.333 & 0.535 & 0.535 \\
\bottomrule\end{tabular}

\end{table}

\textbf{H1 (max+steps beats sum+steps at S64): not supported.} The group means (Table~\ref{tab:hyp}) favour max+steps only because of sum+steps seed 4; the Welch and paired tests do not reject equality ($p=0.374$), and the medians point the other way (Table~\ref{tab:robust}: S64 $0.460$ for sum+steps, $1.765$ for max+steps, with 1/5 and 3/5 failing seeds). The post hoc S16 comparison is not significant either. On the growing-degree family the two systems fail differently (Table~\ref{tab:perseed}): sum+steps gave a non-finite D64 error in 4/5 seeds and an error of order $10^{17}$ in the fifth, whereas max+steps exceeded the threshold in 3/5 seeds with errors that stayed finite (at most $35.55$). No test on D64 is possible with non-finite values.

\textbf{H2 (step supervision helps max): not supported.} In the registered S64 test the means are $0.909$ without and $1.955$ with steps ($p=0.363$): the sign is opposite to H2 and the difference is within seed noise. The exploratory evidence points against H2: without steps the max MPNN has the lowest in-distribution error (S8 median $0.006$ against $0.021$ with steps) and a lower S16 error (paired $p=0.041$ in a post hoc, uncorrected test on five seeds, Table~\ref{tab:hyp}). The medians at S8 and S16 are also lower without steps for sum (S8 $0.018$ vs.\ $0.036$, S16 $0.028$ vs.\ $0.108$). We do not know why: the step loss averages $T$ terms plus a classification loss under the same step budget, so it may simply be optimised less on the final output; we did not run a loss-weight or longer-training control, so this is a statement about this budget only.

\textbf{H3 (MLP worse than every MPNN for $n\ge16$): supported at S16 against max+steps, not supported at S64.} At S16 the MLP is far worse than the reference (MAE $0.671$ vs.\ $0.078$ for max+steps; Welch $p=\rhval{cmp/main/mlp-flat-adjacency/bellman_ford/mae_sparse_n16/welch_p:sci1}$, paired $p=\rhval{cmp/main/mlp-flat-adjacency/bellman_ford/mae_sparse_n16/paired_p:sci1}$); the other MPNNs were not tested against the MLP, but their S16 means lie between $0.03$ and $0.12$ (Table~\ref{tab:main}). The MLP's relative error at 64 nodes (rS64 $0.943$, rD64 $0.926$) is close to that of predicting zero, which is 1 by definition. At S64 its median $1.328$ is \emph{below} that of max+steps ($1.765$), and the test does not separate them ($p=0.535$). The MLP also fits poorly in distribution (S8 $0.239$), so it gives no evidence about extrapolation per se; it may be under-trained at 1000 steps and we did not tune it.

\textbf{H4 (sum degrades more than max under growing degree): direction as hypothesised, untested.} Table~\ref{tab:perseed} lists the two errors that form the per-seed ratio D64/D8; we do not print the quotient. For the registered pair, the step-supervised systems, the D8 errors are of similar size (medians $\rhval{main/mpnn-sum-+-steps/bellman_ford/mae_dense_n8/median:3}$ for sum+steps and $\rhval{main/mpnn-max-+-steps/bellman_ford/mae_dense_n8/median:3}$ for max+steps), whereas the D64 error of sum+steps is non-finite in four seeds and of order $10^{17}$ in the fifth and that of max+steps stays finite (at most $\rhval{main/mpnn-max-+-steps/bellman_ford/mae_dense_n64/max:2}$); the degradation from 8 to 64 nodes is therefore larger for sum in each of the five seeds. No test was registered and none is meaningful on non-finite values, so we report the direction and do not call H4 confirmed. For the final-only pair (post hoc) the per-seed values in Table~\ref{tab:perseed} do not separate the two systems (their D64 errors overlap, with one non-finite seed for sum and two seeds above the threshold for max), and the failure counts at D64 point the other way (1/5 for sum, 2/5 for max). For that pair H4 is inconclusive.

\textbf{Where it fails.} Table~\ref{tab:perseed} shows that every MPNN configuration has at least one failing seed in at least one family. The max MPNN without steps is accurate on three seeds (S64 MAE at most $0.073$) but seeds 0 and 1 fail at both S64 and D64. The sum MPNN without steps has no failing sparse seed, but seed 4 gave a non-finite error at D64. Max+steps fails on seeds 1--3 in both families. Sum+steps fails on one sparse seed and on every dense seed. The non-finite errors come from non-finite network outputs; in the main group they occur only for sum aggregation and only on the dense family at 64 nodes, where a node sums over the most neighbours for the most steps. This may reflect activations growing with degree, but we logged no activations and no run isolates the cause. With five seeds we cannot rank the four MPNN configurations by reliability.
